\documentclass[12pt,a4paper]{report}
\usepackage[utf8]{inputenc}
\usepackage{amsfonts}
\usepackage{setspace}
\usepackage{graphicx}
\usepackage{array}
\usepackage{fancyhdr}
\linespread{1.5}
\usepackage{geometry}
\geometry{
a4paper,
total={210mm,297mm},
left=1.25in,
right=0.75in,
top=0.75in,
bottom=0.75in,
}

\begin{document}
\pagestyle{empty}
\begin{center}
\includegraphics[width=6in,height=1in]{logo.jpg}
\par
\vspace{3pt}
{\large \textbf{(Off-Campus Center, Nitte Deemed to be University, Nitte - 574110)}}
\begin{figure}[hbtp]
\centering
\end{figure}
\\
\textbf{A MINI PROJECT REPORT}
\par
\textbf{on}
\par
\vspace{3pt}
{\Large \textbf{``CAR RENTAL MANAGEMENT SYSTEM''}}
\par
\vspace{8pt}
\textit{\textbf{Submitted by}}
\par
\vspace{6pt}
\textbf{\large Dishanth Shetty}\hspace{3cm}\textbf{\large NNM24CS506}\\
\textbf{\large Akshay Karunakar Naik}\hspace{1cm}\textbf{\large NNM24CS502}\\
\par
\vspace{3pt}
\textit{\textbf{In partial fulfillment of the requirements for the IV semester }}
\par
\vspace{0.5pt}
\Large \textbf{DBMS LABORATORY WITH MINI PROJECT}
\par
\vspace{0.5pt}
\normalsize \centering \textbf{of}
\par
\vspace{0.5pt}
\large \textbf{BACHELOR OF TECHNOLOGY }
\par
\vspace{0.5pt}
\textbf{in}
\par
\vspace{0.5pt}
\large \textbf{COMPUTER SCIENCE \& ENGINEERING}
\par
\vspace{2pt}
\textit{\textbf{Under the Guidance of}}
\par
\vspace{2pt}
\textbf{Ms. Jayapadmini Kanchan}
\par
\vspace{0.5pt}
\textbf{\centering \begin{normalsize}
Assistant Professor Gd-II, Department of CSE
\end{normalsize} }
\par
\vspace{0.5pt}
\normalsize \centering \textbf{at}\\
\begin{figure}[hbtp]
\centering
\includegraphics[width=6in,height=1in]{logo.jpg}
\end{figure}

\par
\vspace{3pt}

\newpage

{\LARGE \textbf{NMAM INSTITUTE OF TECHNOLOGY}}
\par
{\large \textbf{(Deemed to be University)}}
\\
{\large \textbf{Nitte, Karkala}}
\par
\vspace{0.26in}
{\large \textbf{Department of Computer Science \& Engineering}}
\vspace{0.26in}\\
\includegraphics[width=6in,height=1in]{logo.jpg}
\par
\begin{figure}[hbtp]
\centering
\end{figure}
{\Large \textbf{CERTIFICATE}}
\end{center}
\par
\vspace{0.10in}
\setstretch{1.5}
\noindent This is to certify that the \textbf{Mini Project} entitled \textbf{``Car Rental Management System''}  has been carried out by
\textbf{ Dishanth Shetty (NNM24CS506)} and \textbf{Akshay Karunakar Naik  (NNM24CS502)}, the bonafide students of NMAM Institute of Technology in partial fulfillment of the requirements for the IV semester \textbf{DBMS Laboratory with Mini Project (CS2102-1)} of \textbf{Bachelor of Technology in Computer Science \& Engineering} at Nitte Deemed to be University, Nitte during the year 2024 - 25. It is certified that all corrections/suggestions indicated for Internal Assessment have been incorporated in the report deposited in the departmental library. The mini project report has been approved as it satisfies the academic requirements in respect of mini project work.\par
\vspace{0.55in}
\setstretch{1.15}
\begin{tabbing}
-------------------------------------------\hspace{2.0in}--------------------------------------\\
\hspace{0.035in}\textbf{    Ms. Jayapadmini Kanchan}\hspace{2.0in}\textbf{Dr. Jyothi Shetty}\\
\hspace{0.30in}Assistant Professor\hspace{3.0in}Professor \& Head\\
\hspace{0.21in}Dept. of CSE, NMAMIT\hspace{2.4in}Dept. of CSE, NMAMIT\\

\end{tabbing}

\par
\vspace{0.25in}

\newpage
\begin{center}
{\LARGE \textbf{NMAM INSTITUTE OF TECHNOLOGY}}
\par
{\large \textbf{(Deemed to be University)}}
\\
{\large \textbf{Nitte, Karkala}}
\par
\vspace{0.26in}
{\large \textbf{Department of Computer Science \& Engineering}}
\vspace{0.26in}\\
\includegraphics[width=6in,height=1in]{logo.jpg}
\par
\begin{figure}[hbtp]
\centering

\end{figure}
{\Large \textbf{DECLARATION}} 
\end{center}
\par
\vspace{0.10in}
\setstretch{1.5}
\noindent We hereby declare that the entire work embodied in this Mini Project Report titled
\textbf{``Car Rental Management System''} has been carried out by us at NMAM Institute of Technology, Nitte under the supervision of \textbf{Ms. Jayapadmini Kanchan} as the part of the IV semester \textbf{DBMS Laboratory with Mini Project (CS2102-1)} of \textbf{Bachelor of Engineering} in \textbf{Computer Science \& Engineering}. This report has not been submitted to this or any other University.\\
\vspace{0.25in}
\begin{flushright}
\textbf{Dishanth Shetty (NNM24CS506)}\\
\textbf{Akshay Karunakar Naik (NNM24CS502)}\\
NMAMIT, Nitte \\
\end{flushright}

% For Individuals
%\begin{flushright}
%Place: Mangaluru \hspace{7.8cm} \textbf{Student Name 1}\\
%Date : \hspace{11.7cm}4SF11IS001\\
%V Semester, B.E., ISE\\
%SCEM, Mangaluru\\
%\end{flushright}

\newpage
\pagestyle{plain}
\pagenumbering{roman}
\chapter*{Abstract}
\addcontentsline{toc}{chapter}{\numberline{}Abstract}
The Car Rental Management System is a comprehensive and user-friendly software application developed to automate and manage the operations of a car rental service. This system is designed to replace the traditional manual methods of handling rental services, which are often time-consuming, error-prone, and inefficient. By digitizing the process, the system ensures better organization, faster service delivery, and improved customer satisfaction. The core objective of this system is to enable seamless interaction between customers and the rental service provider. Customers can rent vehicles based on availability, and the service provider can manage rental details, customer information, vehicle data, and payment records with greater ease. The system supports the creation and management of bookings, tracks rental durations, monitors the availability of vehicles, and handles payments securely. It also provides the ability to manage multiple branches or locations, ensuring that each vehicle can be tracked according to its assigned area. One of the major advantages of this system is its ability to provide real-time access to accurate information, which helps in making informed business decisions. It minimizes human error, increases data consistency, and reduces the administrative burden on staff. Additionally, the system offers scalability, allowing it to be extended with features like notification alerts, customer feedback, data analytics, and report generation as per future business needs. The Car Rental Management System ultimately leads to enhanced operational efficiency, better resource utilization, and a professional rental experience for customers. It serves as a reliable solution for car rental companies aiming to modernize their services, stay competitive, and meet the growing expectations of today’s tech-savvy users.
\chapter*{Acknowledgement}
\addcontentsline{toc}{chapter}{\numberline{}Acknowledgement}
We dedicate this page to acknowledging and thanking those responsible for shaping the project. Without their guidance and help, the experience while constructing the dissertation would not have been so smooth and efficient.
\par
\vspace{0.15in}
\noindent We sincerely thank  \textbf{Dr. Niranjan N. Chiplunkar}, Principal, NMAM Institute of Technology, Nitte who have always been a great source of inspiration.\vspace{0.15in}\\
\noindent We convey our sincere regards to \textbf{Dr. I. R. Mithanthaya}, Vice Principal, NMAM Institute of Technology, Nitte for encouraging us during the project tenure. \vspace{0.15in}\\
\noindent We express our sincere gratitude to \textbf{Dr. Jyothi Shetty}, Head of the Department, Computer Science and Engineering, NMAM Institute of Technology, Nitte for her valuable support, guidance, encouragement through the project.
\vspace{0.15in}\\
\noindent We take this opportunity to express our deep sense of gratitude to our guide \textbf{Ms. Jayapadmini Kanchan}, Assistant Professor Gd-II, Department of Computer Science \& Engineering, NMAMIT, Nitte for her whole-hearted support and valuable guidance she has given us during the project part.
\par
\vspace{0.15in}
\noindent Finally, yet importantly, We express our heartfelt thanks to our family \& friends for their wishes and encouragement throughout the work.
\\
\\
%\begin{flushright}
%\textbf{Student Name 1}\\
%\textbf{Student Name 2}\\
%\end{flushright}

% For Individuals
\begin{flushright}
\begin{tabular}{p{10cm} p{10cm}}
\textbf{Dishanth Shetty} & \textbf{Akshay Karunakar Naik} \\
NNM24CS506 & NNM24CS502 \\
IV Sem, B.Tech. CSE & IV Sem, B.Tech. CSE \\
NMAMIT, Nitte & NMAMIT, Nitte \\
\end{tabular}
\end{flushright}


\setstretch{1.4}
\renewcommand{\contentsname}{Table of Contents}
\tableofcontents
\addcontentsline{toc}{chapter}{\numberline{}Table of Contents}
\listoffigures
\addcontentsline{toc}{chapter}{\numberline{}List of Figures}


\pagestyle{fancy}
\fancyhf{}
\lhead{\fontsize{10}{12} \selectfont Car Rental Management System}
\rhead{\fontsize{10}{12} \selectfont Chapter \thechapter}
\lfoot{\fontsize{10}{12} \selectfont Department of Computer Science \& Engineering, NMAMIT, Nitte}
\rfoot{\fontsize{10}{12} \selectfont Page \thepage}
\renewcommand{\headrulewidth}{0.5pt}
\renewcommand{\footrulewidth}{0.5pt}

\setstretch{1.5}

\chapter{Introduction}
\pagenumbering{arabic}
\par
In today’s fast-paced and technology-driven world, the demand for convenient and efficient transportation services is on the rise. Car rental services have become an essential part of urban life, offering individuals and businesses flexible mobility options without the long-term commitment of vehicle ownership. Managing such services manually often leads to challenges such as delayed bookings, errors in record-keeping, and inefficient utilization of resources. 
\\
\para\noindent The Car Rental Management System is a software-based solution developed to automate and streamline the operations of a car rental service. It aims to improve the overall efficiency of the system by digitally managing customer records, vehicle availability, bookings, rental transactions, and branch-level coordination. This system reduces the dependency on paperwork and manual processes, ensuring accuracy, consistency, and real-time access to essential data.
\par
\section{Purpose}
The primary purpose of this system is to automate and optimize the process of renting vehicles to customers in a structured, reliable, and efficient manner. This system is developed to address the limitations of traditional manual methods, such as errors in record-keeping, delays in processing bookings, and inefficient management of vehicle availability. By implementing this system, rental service providers can maintain a centralized database that manages customer details, vehicle records, rental transactions, and payment history with ease. It ensures real-time tracking of vehicle availability across different locations or branches, reduces administrative workload, and improves overall operational efficiency.


\section{Scope}
This system is developed to provide an efficient and streamlined approach to managing the operations of a car rental business. The system is designed to handle various aspects of the rental process, including customer management, vehicle tracking, rental bookings, payments, and location-based availability of vehicles. It serves both administrative users and customers by offering a centralized platform to monitor and manage all transactions.
The scope of this project covers the complete life cycle of a car rental transaction — from the moment a customer initiates a booking to the final payment and return of the vehicle. The system supports multiple customers and vehicles simultaneously, with features that ensure each car is linked to a specific location and rental record. Administrators can access rental history, check vehicle status, assign vehicles to different branches, and handle customer queries efficiently.
\section{Overview}
This system is a software-based solution designed to automate and simplify the day-to-day operations of a car rental service. It provides an organized platform for managing customers, vehicles, rental transactions, payments, and location details. This system replaces the traditional manual process, which is often inefficient, time-consuming, and prone to errors. With this system, customers can easily rent vehicles based on availability, while the rental service provider can manage and monitor all operations through a centralized interface. It maintains records of customer details, car information, rental dates and payment statuses. Additionally, it ensures that each car is linked to a specific location, making it easier to track and manage resources across multiple branches.

\chapter{Requirements Specification}
\section{Hardware Specification}
\begin{itemize}
\item Processor : Intel Core i3 or higher
\item RAM : Minimum 4 GB (8 GB recommended)
\item Hard Disk : Minimum 250 GB
\item Input Device : Standard keyboard and Mouse
\item Output Device : Monitor
\end{itemize}
\section{Software Specification}
\begin{itemize}
\item Database : MySQL Workbench 8.0 CE
\item Markup Language :  HTML
\item Scripting Language: JSP, Servlet, Javascript
\item IDE : Eclipse IDE for Enterprise Java and Web Developers - 2023 - 12
\end{itemize}

\chapter{System Design}
\section{ER Diagram}
\begin{figure}[hbtp]
\centering
\includegraphics[width=4in,height=3in]{Car_Rental_ER.png}
\caption{ER Diagram for Car Rental Management System}
\end{figure} 
\para\noindent Customer will interact with the system through the application interface to rent a car by providing personal details and selecting rental preferences. Once the booking is initiated, the system creates a rental record and links it with the selected car and customer. The admin monitors the rental activity and verifies car availability. Based on the location, car availability is fetched from the database. Once the rental is confirmed, the admin ensures payment details are recorded, and the rental is marked with the appropriate status. The main functionality, managing rental processes and ensuring car availability across locations, is handled by the admin.
\section{Mapping From ER Diagram to Schema Diagram}

\begin{enumerate}
    \item \textbf{Mapping of Regular (Strong) Entity Types:} \\
    Each strong entity becomes a table with its attributes; the primary key stays the same.
    
    \item \textbf{Mapping of Weak Entity Types:} \\
    A weak entity becomes a table with a foreign key from the strong entity and its partial key.

    \item \textbf{Mapping of Binary 1:1 Relationship Types:} \\
    For a 1:1 relationship, add the primary key of one entity as a foreign key in the other.

    \item \textbf{Mapping of Binary 1:N Relationship Types:} \\
    For a 1:N relationship, add the primary key of the "one" side as a foreign key in the "many" side.

    \item \textbf{Mapping of Binary M:N Relationship Types:} \\
    For an M:N relationship, create a new table with foreign keys from both entities and attributes.

    \item \textbf{Mapping of Multi-Valued Attributes:} \\
    For multi-valued attributes, create a new table with the attribute and a foreign key to the entity.

    \item \textbf{Mapping of N-ary Relationship Types:} \\
    For n-ary relationships, create a new table with foreign keys from all participating entities.
\end{enumerate}


\section{Assumptions}
\begin{itemize}
	\item One customer can rent multiple cars.
        \item One rental is linked to only one customer.
        \item One car can be rented multiple times, but only once at a time.
        \item Each rental must be associated with one payment.
        \item One car belongs to only one branch.
        \item Not all customers may have active rentals.
        \item One admin can handle multiple rental transactions.
        \item All cars may not be available for rent at all times.

	
\end{itemize}
\section{Schema Diagram }
A Data Flow Diagram (DFD) represents the flow of data within the Car Rental Management System. It shows how data moves between customers, admins, and system components like rental, payment, car, and location. Data flows from external entities (like customers) into internal processes (such as booking a car or making payments) and is stored in data stores like rental records and customer databases. The DFD focuses on data movement and transformation, not on process order or timing. Unlike a flowchart, it emphasizes where data is stored, what data is processed, and how it travels through the system.
\begin{figure}[hbtp]
\centering
\includegraphics[width=6.6in,height=4.6in]{Car_Rental_Schema.png}
\caption{Schema Diagram for Car Rental Management System}
\end{figure}
\chapter{Implementation}
\section{Pseudocode Used}
This System consists of several modules that align with the entities defined in the ER model. These modules include \textbf{Customer}, \textbf{Car}, \textbf{Rental}, \textbf{Payment}, \textbf{Location}, and \textbf{Admin}. Each module performs specific tasks to ensure the smooth functioning of the system.

\section*{Customer Module}
This module allows new customers to register by entering their personal details such as name, address, phone number, and email. A unique Customer ID (CUST\_ID) is generated for each new customer and stored in the \texttt{CUSTOMER} table.

\textbf{Pseudocode:}
\begin{itemize}
    \item Start
    \item Enter customer details (name, address, phone, email)
    \item Generate \texttt{CUST\_ID}
    \item Insert into \texttt{CUSTOMER} table
    \item End
\end{itemize}

\section*{Car Module}
This module handles the management of cars available for rent. Each car is identified by a unique registration number and is associated with a location. When a customer makes a booking request, the system checks the availability of cars based on the selected location and model.

\textbf{Pseudocode:}
\begin{itemize}
    \item Start
    \item Enter location and car model
    \item Check available cars from \texttt{CAR} table
    \item If available, assign car and update status to ``Booked''
    \item End
\end{itemize}

\section*{Rental Module}
This module processes car booking requests. It collects necessary details such as the Customer ID, Car Registration Number, Start Date, and End Date. A unique Rental ID is created and stored in the \texttt{RENTAL} table.

\textbf{Pseudocode:}
\begin{itemize}
    \item Start
    \item Enter booking details (CUST\_ID, C\_REG\_NO, Start Date, End Date)
    \item Generate \texttt{RENTAL\_ID}
    \item Insert into \texttt{RENTAL} table
    \item End
\end{itemize}

\section*{Payment Module}
This module is responsible for handling payment transactions. It takes the Rental ID and payment method as input, generates a Payment ID, records the transaction date and status, and stores the data in the \texttt{PAYMENT} table.
\\
\\
\textbf{Pseudocode:}
\begin{itemize}
    \item Start
    \item Enter payment details (RENTAL\_ID, payment method)
    \item Generate \texttt{PAY\_ID}
    \item Record date and status
    \item Insert into \texttt{PAYMENT} table
    \item End
\end{itemize}

\section*{Car Return Module}
Once the rental period is over, this module updates the rental status to ``Completed'' and sets the corresponding car status back to ``Available'' in the database.

\textbf{Pseudocode:}
\begin{itemize}
    \item Start
    \item Enter \texttt{RENTAL\_ID}
    \item Update \texttt{RENTAL} table: set \texttt{R\_STATUS} = ``Completed''
    \item Find \texttt{C\_REG\_NO} and update \texttt{CAR} status to ``Available''
    \item End
\end{itemize}

\para\noindent This modular implementation ensures efficient and organized management of the Car Rental Management System. Each module is designed to handle specific operations while working together to provide a seamless experience for users and administrators.
\\
\\
\\

\section{Tables Used}
\par This car rental database neatly organizes key details into five main sections. CUSTOMER simply lists renter information, while LOCATION shows where cars are available. The CAR section tracks each vehicle's details and its current status and rental place. RENTAL agreements record important dates, who rented the car, and which car they took. Lastly, PAYMENT tracks money details for each rental. Unique IDs (#) help link everything together smoothly, making it easy to manage the car rental process.\\


\begin{itemize}
  \item \textbf{CUSTOMER} ( \underline{custId}: Int, custName: Varchar, cAddress: Varchar, custPhone: Varchar, custEmail: Varchar )
  
  \item \textbf{LOCATION} ( \underline{locId}: Int, locName: Varchar, locCity: Varchar, locState: Varchar, locZip: Varchar )
  
  \item \textbf{CAR} ( \underline{cRegNo}: Varchar, carModel: Varchar, carYear: Int, carColor: Varchar, carStatus: Varchar, locId: Int, custId: Int )
  
  \item \textbf{RENTAL} ( \underline{rentalId}: Int, rName: Varchar, startDate: Date, endDate: Date, rStatus: Varchar, custId: Int, cRegNo: Varchar )
  
  \item \textbf{PAYMENT} ( \underline{payId}: Int, dateTime: Datetime, payMethod: Varchar, payStatus: Varchar, rentalId: Int )
\end{itemize}

\chapter{Results and Disscussion}
\par The Figure 5.1 represents the home landing page of the \textbf{Car Rental Management System}, which provides a user-friendly interface featuring prominent search options for vehicle selection. It typically includes featured vehicles, promotional banners, and clear calls to action for seamless booking and navigation.
\begin{figure}[hbtp]
    \centering
    \includegraphics[width=0.8\textwidth, keepaspectratio]{HomePage.jpeg}
    \caption{Home Landing Page for Car Rental Management System}
    \label{fig:landing_page}
\end{figure}

\\
\\
\vspace{0.15in}

\par
\para\noindent Figure 5.2 represents the Customer Registration page of the \textbf{Car Rental Management System}. This interface is designed to capture essential user information required for account creation and future interactions with the platform. The form includes clearly labeled input fields for key details such as the customer’s full name, address, phone number, and email address. These fields ensure that the system collects accurate and necessary data to facilitate bookings, communication, and customer service.

\begin{figure}[hbtp]
    \centering
    \includegraphics[width=0.8\textwidth, keepaspectratio]{Customer-Registration.jpeg}
    \caption{Customer Registration}
    \label{fig:customer_register}
\end{figure}
\para\noindent Figure 5.3 represents the Customer Login page of the \textbf{Car Rental Management System}. This page enables registered users to securely access their accounts by providing their login credentials. The interface includes input fields for essential details such as Customer ID (CustID) and email address, which are required to authenticate the user. The layout is designed to be simple and intuitive, allowing users to quickly log in without confusion or delay.In addition to the input fields, the page may feature buttons like “Login” or “Sign In” to initiate the authentication process.

\begin{figure}[hbtp]
    \centering
    \includegraphics[width=0.8\textwidth, keepaspectratio]{Customer-Login.jpeg}
    \caption{Customer Login}
    \label{fig:customer_login}
\end{figure}
\\
\para\noindent Figure 5.4 represents the Customer Dashboard page of the \textbf{Car Rental Management System}. This page provides users with a centralized and comprehensive overview of their interactions within the platform. The dashboard displays important information such as the customer’s booking history, active rentals, and profile details. This allows users to keep track of their past and current rentals with ease. The layout is designed to be intuitive and user-friendly, ensuring that customers can navigate through different sections effortlessly.
\\
\\
\begin{figure}[hbtp]
    \centering
    \includegraphics[width=0.8\textwidth, keepaspectratio]{Customer-Dashboard.jpeg}
    \caption{Customer Dashboard}
    \label{fig:customer_dashboard}
\end{figure}

\para\noindent Figure 5.5 represents the Admin Login page, which includes input fields for the essential login details such as the Admin ID and Email.

\begin{figure}[hbtp]
    \centering
    \includegraphics[width=0.8\textwidth, keepaspectratio]{Admin-Login.jpeg}
    \caption{Admin Login}
    \label{fig:admin_login}
\end{figure}


\\
\para\noindent Figure 5.6 represents the Admin Dashboard page of the \textbf{Car Rental Management System}. This dashboard serves as the central control panel for administrators, offering access to a wide range of backend operations essential for managing the platform.
\begin{figure}[hbtp]
    \centering
    \includegraphics[width=0.8\textwidth, keepaspectratio]{Admin-Dashboard.jpeg}
    \caption{Admin Dashboard}
    \label{fig:admin_dashboard}
\end{figure}

\para\noindent Figure 5.7 represents the Manage Cars page of the \textbf{Car Rental Management System}. This interface is specifically designed for administrators to efficiently oversee and manage the platform’s vehicle inventory. The page displays a comprehensive list of all available vehicles along with key details such as car model, registration number, availability status, rental price, and more.

\\

\begin{figure}[hbtp]
    \centering
    \includegraphics[width=0.8\textwidth, keepaspectratio]{Admin-Manage-Cars.jpeg}
    \caption{Manage Cars - Admin Panel}
    \label{fig:manage_cars}
\end{figure}

\para\noindent Figure 5.8 represents the Manage Rentals page of the \textbf{Car Rental Management System}. This section of the admin interface is dedicated to overseeing and managing all rental transactions processed through the platform. The page provides administrators with a detailed overview of both ongoing and completed rentals, including information such as rental ID, customer name, vehicle details, rental duration, status, and payment information. 

\\

\begin{figure}[hbtp]
    \centering
    \includegraphics[width=0.8\textwidth, keepaspectratio]{Admin-Manage-Rentals.jpeg}
    \caption{Manage Rentals - Admin Panel}
    \label{fig:manage_rentals}
\end{figure}

\\

\para\noindent Figure 5.9 represents the Manage Customers page of the Car Rental Management System. This administrative interface is designed to provide the Admin with full control over user management by displaying a comprehensive list of all registered customers. The page showcases key customer details such as full name, contact information, email address, and registration status. It includes essential functionalities that allow the administrator to view, update, or delete customer records as needed, ensuring that the system's user database remains accurate and up to date.

\\

\begin{figure}[hbtp]
    \centering
    \includegraphics[width=0.8\textwidth, keepaspectratio]{Admin-Customer-Management.jpeg}
    \caption{Manage Customers - Admin Panel}
    \label{fig:manage_customers}
\end{figure}


\para\noindent Figure 5.10 represents the Manage Locations page of the \textbf{Car Rental Management System}. This page allows the administrator to view a comprehensive list of all available pick-up and drop-off locations associated with the car rental service. The interface presents an organized and clear display of location details, including the location name, address, and other relevant information. This helps the admin easily reference and verify the locations used within the system

\\

\begin{figure}[hbtp]
    \centering
    \includegraphics[width=0.8\textwidth, keepaspectratio]{Admin-Location-Management.jpeg}
    \caption{Manage Locations - Admin Panel}
    \label{fig:manage_locations}
\end{figure}


\chapter{Conclusion and Future work}
\p The development of the \textbf{Car Rental Management System} has successfully streamlined the rental process by integrating key functionalities such as vehicle inventory management, customer profiling, reservation handling, and secure payment processing. By leveraging a robust relational database schema, the system ensures efficient data organization, facilitating real-time tracking of vehicle availability and comprehensive management of rental agreements. This structured approach not only enhances operational efficiency but also improves data accuracy and customer satisfaction. The implementation demonstrates a significant advancement over traditional manual processes, offering a scalable and user-friendly solution that meets the dynamic needs of the car rental industry. 
\\
\\
However, there is potential for \textbf{future enhancements}, such as the development of a mobile application to increase customer accessibility, the introduction of advanced analytics and reporting tools to provide deeper insights into rental trends, and the integration of GPS tracking systems for better fleet management. Additionally, incorporating new payment methods, integrating AI-driven recommendations, and implementing a customer feedback system could further improve the platform’s functionality, offering more personalized services and driving continuous business growth. These future advancements would ensure that the system remains adaptive, scalable, and aligned with the evolving needs of the car rental industry.
\newpage
\pagestyle{plain}
\renewcommand{\bibname}{References}
\bibliography{references}
\addcontentsline{toc}{chapter}{References}
\begin{thebibliography}{35}

\bibitem{mor}
Coronel, C., Morris, S., and Rob, P. Database Systems: Design, Implementation, and Management. Cengage Learning, 2012.

\bibitem{jsp}
Hanna, P. JSP: The Complete Reference. 1st ed., McGraw-Hill Education, 2001.

\bibitem{servlet}
Hunter, J., and Crawford, W. Java Servlet Programming. 2nd ed., O’Reilly Media, 2001.

\bibitem{jdbc}
Fisher, C., Ellis, J., and Rees, J. JDBC API Tutorial and Reference. 3rd ed., Addison-Wesley, 2003.

\bibitem{eclipse}
Eclipse Foundation. Eclipse IDE Documentation. Available at: \url{https://www.eclipse.org}

\bibitem{mysql}
Oracle Corporation. MySQL Workbench Manual. Available at: \url{https://dev.mysql.com/doc/workbench/en/}

\end{thebibliography}


\end{document}